\documentclass[10pt,letterpaper]{article}
\usepackage[utf8]{inputenc}
\usepackage[intlimits]{amsmath}
\usepackage{amsfonts}
\usepackage{amssymb}
\usepackage{ragged2e}
\usepackage[letterpaper, margin=0.5in]{geometry}
\usepackage{graphicx}
\usepackage{cancel}
\usepackage{mathtools}
\usepackage{tabularx}
\usepackage{arydshln}
\usepackage{tensor}
\usepackage{array}
\usepackage{xcolor}
\usepackage[boxed]{algorithm}
\usepackage[noend]{algpseudocode}
\usepackage{listings}
\usepackage{textcomp}
% \usepackage[pdf,tmpdir,singlefile]{graphviz}
\usepackage{mathrsfs}
\usepackage{bbm}
\usepackage{tikz}
\usepackage{tikz-cd}
\usepackage{enumitem}
\usepackage{arydshln}
\usepackage{relsize}
\usepackage{multicol}
\usepackage{scalerel}
\usepackage{upgreek}
\usepackage{wasysym}

\usetikzlibrary{bayesnet}
\setlist{noitemsep}

%%%%%%%%%%%%%%%%%%%%%%%%%%%%%
% Formatting commands
%%%%%%%%%%%%%%%%%%%%%%%%%%%%%
\newcommand{\n}{\hfill\break}
\newcommand{\nn}{\vspace{0.5\baselineskip}\n}
\newcommand{\up}{\vspace{-\baselineskip}}
\newcommand{\hangblock}[2]{\par\noindent\settowidth{\hangindent}{\textbf{#1: }}\textbf{#1: }\nolinebreak#2}
\newcommand{\lemma}[1]{\hangblock{Lemma}{#1}}
\newcommand{\defn}[1]{\hangblock{Defn}{#1}}
\newcommand{\thm}[1]{\hangblock{Thm}{#1}}
\newcommand{\cor}[1]{\hangblock{Cor}{#1}}
\newcommand{\prop}[1]{\hangblock{Prop}{#1}}
\newcommand{\ex}[1]{\hangblock{Ex}{#1}}
\newcommand{\exer}[1]{\hangblock{Exer}{#1}}
\newcommand{\fact}[1]{\hangblock{Fact}{#1}}
\newcommand{\remark}[1]{\hangblock{Remark}{#1}}
\newcommand{\proven}{\;$\square$\n}
\newcommand{\problem}[1]{\par\noindent{#1}\n}
\newcommand{\problempart}[2]{\par\noindent\indent{}\settowidth{\hangindent}{\textbf{(#1)} \indent{}}\textbf{(#1) }\nolinebreak#2\n}
\newcommand{\ptxt}[1]{\textrm{\textnormal{#1}}}
\newcommand{\inlineeq}[1]{\centerline{$\displaystyle #1$}}
\newcommand{\pageline}{\up\par\noindent\rule{\textwidth}{0.1pt}}

%%%%%%%%%%%%%%%%%%%%%%%%%%%%%
% Math commands
%%%%%%%%%%%%%%%%%%%%%%%%%%%%%
% Set Theory
\newcommand{\card}[1]{\left|#1\right|}
\newcommand{\set}[1]{\left\{#1\right\}}
\newcommand{\setmid}{\;\middle|\;}
\newcommand{\ps}[1]{\mathcal{P}\left(#1\right)}
\newcommand{\pfinite}[1]{\mathcal{P}^{\ptxt{finite}}\left(#1\right)}
\newcommand{\naturals}{\mathbb{N}}
\newcommand{\N}{\naturals}
\newcommand{\integers}{\mathbb{Z}}
\newcommand{\Z}{\integers}
\newcommand{\rationals}{\mathbb{Q}}
\newcommand{\Q}{\rationals}
\newcommand{\reals}{\mathbb{R}}
\newcommand{\R}{\reals}
\newcommand{\complex}{\mathbb{C}}
\newcommand{\C}{\complex}
\newcommand{\halfPlane}{\mathbb{H}}
\let\HSym\H
\let\H\relax
\newcommand{\H}{\halfPlane}
\newcommand{\comp}{^{\complement}}
\DeclareMathOperator{\Hom}{Hom}
\newcommand{\Ind}{\mathbbm{1}}
\newcommand{\cut}{\setminus}
\DeclareMathOperator{\elem}{elem}

% Differentiable Manifolds
\newcommand{\RP}{\mathbb{RP}}
\newcommand{\CP}{\mathbb{RP}}
\newcommand{\osubset}{\overset{\mathclap{\scalebox{0.5}{\ptxt{open}}}}{\subset}}
\newcommand{\osubseteq}{\overset{\mathclap{\scalebox{0.5}{\ptxt{open}}}}{\subseteq}}
\newcommand{\osupset}{\overset{\mathclap{\scalebox{0.5}{\ptxt{open}}}}{\supset}}
\newcommand{\osupseteq}{\overset{\mathclap{\scalebox{0.5}{\ptxt{open}}}}{\supseteq}}
\newcommand{\pdat}[3]{\left.\pd{#1}{#2}\right|_{#3}}
\DeclareMathOperator{\so}{so}
\DeclareMathOperator{\codim}{codim}
\DeclareMathOperator{\Diff}{Diff}
\let\dSym\d
\let\d\relax
\newcommand{\d}{\partial}
\DeclareMathOperator{\gl}{gl}

% Graph Theory
\let\deg\relax
\DeclareMathOperator{\deg}{deg}
\newcommand{\degp}{\ptxt{deg}^{+}}
\newcommand{\degn}{\ptxt{deg}^{-}}
\newcommand{\precdot}{\mathrel{\ooalign{$\prec$\cr\hidewidth\hbox{$\cdot\mkern0.5mu$}\cr}}}
\newcommand{\succdot}{\mathrel{\ooalign{$\cdot\mkern0.5mu$\cr\hidewidth\hbox{$\succ$}\cr\phantom{$\succ$}}}}
\DeclareMathOperator{\cl}{cl}
\DeclareMathOperator{\affdim}{affdim}

% Probability
\newcommand{\parSymbol}{\P}
\newcommand{\Prob}{\mathbb{P}}
\renewcommand{\P}{\Prob}
\newcommand{\Avg}{\mathbb{E}}
\newcommand{\E}{\Avg}
\DeclareMathOperator{\Var}{Var}
\DeclareMathOperator{\cov}{cov}
\DeclareMathOperator{\Unif}{Unif}
\DeclareMathOperator{\Binom}{Binom}
\newcommand{\CI}{\mathrel{\text{\scalebox{1.07}{$\perp\mkern-10mu\perp$}}}}

% Standard Math
\newcommand{\inv}{^{-1}}
\newcommand{\abs}[1]{\left|#1\right|}
\newcommand{\ceil}[1]{\left\lceil{}#1\right\rceil{}}
\newcommand{\floor}[1]{\left\lfloor{}#1\right\rfloor{}}
\newcommand{\conj}[1]{\overline{#1}}
\newcommand{\of}{\circ}
\newcommand{\tri}{\triangle}
\newcommand{\inj}{\hookrightarrow}
\newcommand{\surj}{\twoheadrightarrow}
\newcommand{\ndiv}{\nmid}
\renewcommand{\epsilon}{\varepsilon}
\newcommand{\divides}{\mid}
\newcommand{\ndivides}{\nmid}
\DeclareMathOperator{\lcm}{lcm}
\DeclareMathOperator{\sgn}{sgn}
\newcommand{\map}[4]{\!\!\!\begin{array}[t]{rcl}#1 & \!\!\!\!\to & \!\!\!\!#2\\ {}#3 & \!\!\!\!\mapsto & \!\!\!\!#4\end{array}}
\newcommand{\bigsum}[2]{\smashoperator[lr]{\sum_{\scalebox{#1}{$#2$}}}}
\DeclareMathOperator{\gcf}{gcf}
\newcommand{\restr}[2]{\left.#1\right|_{#2}}

% Linear Algebra
\newcommand{\Id}{\textrm{\textnormal{Id}}}
\newcommand{\im}{\textrm{\textnormal{im}}}
\newcommand{\norm}[1]{\abs{\abs{#1}}}
\newcommand{\tpose}{^{T}\!}
\newcommand{\iprod}[1]{\left<#1\right>}
\newcommand{\giprod}{\iprod{\;\,,\;}}
\DeclareMathOperator{\tr}{tr}
\DeclareMathOperator{\trace}{tr}
\newcommand{\chgBasMat}[3]{\!\!\tensor*[_{#1}]{\left[#2\right]}{_{#3}}}
\newcommand{\vecBas}[2]{\tensor*[]{\left[#1\right]}{_{#2}}}
\DeclareMathOperator{\GL}{GL}
\DeclareMathOperator{\Mat}{Mat}
\DeclareMathOperator{\vspan}{span}
\DeclareMathOperator{\rank}{rank}
\newcommand{\V}[1]{\vec{#1}}
\DeclareMathOperator{\proj}{proj}
\DeclareMathOperator{\compProj}{comp}
\DeclareMathOperator{\row}{row}
\newcommand{\smallPMatrix}[1]{\paren{\begin{smallmatrix}#1\end{smallmatrix}}}
\newcommand{\smallBMatrix}[1]{\brack{\begin{smallmatrix}#1\end{smallmatrix}}}
\newcommand{\pmat}[1]{\begin{pmatrix}#1\end{pmatrix}}
\newcommand{\bmat}[1]{\begin{bmatrix}#1\end{bmatrix}}
\newcommand{\dual}{^{*}}
\newcommand{\pinv}{^{\dagger}}
\newcommand{\horizontalMatrixLine}{\ptxt{\rotatebox[origin=c]{-90}{$|$}}}
\DeclareMathOperator{\range}{range}
\DeclareMathOperator{\Symm}{Symm}
\DeclareMathOperator{\SU}{SU}

% Multilinear Algebra
\let\Lsym\L
\let\L\relax
\DeclareMathOperator{\L}{\mathscr{L}}
\DeclareMathOperator{\A}{\mathcal{A}}
\DeclareMathOperator{\Alt}{Alt}
\DeclareMathOperator{\Sym}{Sym}
\newcommand{\ot}{\otimes}
\newcommand{\ox}{\otimes}
\DeclareMathOperator{\asc}{asc}
\DeclareMathOperator{\asSet}{set}
\DeclareMathOperator{\sort}{sort}
\DeclareMathOperator{\ringA}{\mathring{A}}

% Topology
\newcommand{\closure}[1]{\overline{#1}}
\newcommand{\uball}{\mathcal{U}}
\DeclareMathOperator{\Int}{Int}
\DeclareMathOperator{\Ext}{Ext}
\DeclareMathOperator{\Bd}{Bd}
\DeclareMathOperator{\rInt}{rInt}
\DeclareMathOperator{\ch}{ch}
\DeclareMathOperator{\ah}{ah}
\newcommand{\LargerTau}{\mathlarger{\mathlarger{\mathlarger{\mathlarger{\tau}}}}}
\newcommand{\Tau}{\mathcal{T}}

% Analysis
\DeclareMathOperator{\Graph}{Graph}
\DeclareMathOperator{\epi}{epi}
\DeclareMathOperator{\hypo}{hypo}
\DeclareMathOperator{\supp}{supp}
\newcommand{\lint}[2]{\underset{#1}{\overset{#2}{{\color{black}\underline{{\color{white}\overline{{\color{black}\int}}\color{black}}}}}}}
\newcommand{\uint}[2]{\underset{#1}{\overset{#2}{{\color{white}\underline{{\color{black}\overline{{\color{black}\int}}\color{black}}}}}}}
\newcommand{\alignint}[2]{\underset{#1}{\overset{#2}{{\color{white}\underline{{\color{white}\overline{{\color{black}\int}}\color{black}}}}}}}
\newcommand{\extint}{\ptxt{ext}\int}
\newcommand{\extalignint}[2]{\ptxt{ext}\alignint{#1}{#2}}
\newcommand{\conv}{\ast}
\newcommand{\pd}[2]{\frac{\partial{}#1}{\partial{}#2}}
\newcommand{\del}{\nabla}
\DeclareMathOperator{\grad}{grad}
\DeclareMathOperator{\curl}{curl}
\let\div\relax
\DeclareMathOperator{\div}{div}
\DeclareMathOperator{\vol}{vol}

% Complex Analysis
\let\Re\relax
\DeclareMathOperator{\Re}{Re}
\let\Im\relax
\DeclareMathOperator{\Im}{Im}
\DeclareMathOperator{\Res}{Res}

% Abstract Algebra
\DeclareMathOperator{\ord}{ord}
\newcommand{\generated}[1]{\left<#1\right>}
\newcommand{\cycle}[1]{\smallPMatrix{#1}}
\newcommand{\id}{\ptxt{id}}
\newcommand{\iso}{\cong}
\DeclareMathOperator{\Aut}{Aut}
\DeclareMathOperator{\SL}{SL}
\DeclareMathOperator{\op}{op}
\newcommand{\isom}[4]{\!\!\!\begin{array}[t]{rcl}#1 & \!\!\!\!\overset{\sim}{\to} & \!\!\!\!#2\\ #3 & \!\!\!\!\mapsto & \!\!\!\!#4\end{array}}
\newcommand{\F}{\mathbb{F}}
\newcommand{\acts}{\;\reflectbox{\rotatebox[origin=c]{-90}{$\circlearrowright$}}\;}
\newcommand{\disjunion}{\mathrel{\text{\scalebox{1.07}{$\perp\mkern-10mu\perp$}}}}
\DeclareMathOperator{\SO}{SO}
\DeclareMathOperator{\stab}{stab}
\DeclareMathOperator{\nullity}{nullity}
\DeclareMathOperator{\Perm}{Perm}
\DeclareMathOperator{\nsubgp}{\vartriangleleft}
\DeclareMathOperator{\notnsubgp}{\ntriangleleft}
\newcommand{\presentation}[2]{\left<#1\;\middle|\;#2\right>}
\DeclareMathOperator{\Char}{char}
\DeclareMathOperator{\fchar}{char}
\DeclareMathOperator{\triv}{triv}
\DeclareMathOperator{\reg}{reg}
\DeclareMathOperator{\std}{std}
\DeclareMathOperator{\Func}{Func}
\DeclareMathOperator{\End}{End}

% Convex Optimization
\let\sectionSymbol\S
\let\S\relax
\newcommand{\S}{\mathbb{S}}
\DeclareMathOperator{\dist}{dist}
\DeclareMathOperator{\dom}{dom}
\DeclareMathOperator{\diag}{diag}
\DeclareMathOperator{\ones}{\mathbbm{1}}
\newcommand{\minimizeOver}[3]{\begin{array}{rl}\underset{#1}{\ptxt{minimize}} & #2\\ \ptxt{subject to} & #3\end{array}}
\newcommand{\maximizeOver}[3]{\begin{array}{rl}\underset{#1}{\ptxt{maximize}} & #2\\ \ptxt{subject to} & #3\end{array}}
\newcommand{\minimizationProblem}[2]{\minimizeOver{}{#1}{#2}}
\newcommand{\maximizationProblem}[2]{\maximizeOver{}{#1}{#2}}
\newcommand{\minimizeOverUnconstrained}[2]{\begin{array}{rl}\underset{#1}{\ptxt{minimize}} & #2\end{array}}
\newcommand{\maximizeOverUnconstrained}[2]{\begin{array}{rl}\underset{#1}{\ptxt{maximize}} & #2\end{array}}
\newcommand{\minimizationUnconstrained}[1]{\minimizeOverUnconstrained{}{#1}}
\newcommand{\maximizationUnconstrained}[1]{\maximizeOverUnconstrained{}{#1}}
\DeclareMathOperator{\argmin}{argmin}
\DeclareMathOperator{\argmax}{argmax}

% Proofs
\newcommand{\st}{s.t.}
\newcommand{\unique}{!}
\newcommand{\iffdef}{\overset{\ptxt{def}}{\Leftrightarrow}}
\newcommand{\eqdef}{\overset{\ptxt{def}}{=}}
\newcommand{\eqVertical}{\rotatebox[origin=c]{90}{=}}
\newcommand{\mapsfrom}{\mathrel{\reflectbox{\ensuremath{\mapsto}}}}
\newcommand{\mapsdown}{\rotatebox[origin=c]{-90}{$\mapsto$}\mkern2mu}
\newcommand{\mapsup}{\rotatebox[origin=c]{90}{$\mapsto$}\mkern2mu}
\newcommand{\from}{\!\mathrel{\reflectbox{\ensuremath{\to}}}}
\newcommand{\labeledeq}[1]{\overset{\mathclap{\ptxt{#1}}}{=}}

% Brackets
\newcommand{\paren}[1]{\left(#1\right)}
\renewcommand{\brack}[1]{\left[#1\right]}
\renewcommand{\brace}[1]{\left\{#1\right\}}
\newcommand{\ang}[1]{\left<#1\right>}

% Algorithms
\algrenewcommand{\algorithmiccomment}[1]{\hskip 1em \texttt{// #1}}
\algrenewcommand\algorithmicrequire{\textbf{Input:}}
\algrenewcommand\algorithmicensure{\textbf{Output:}}
\newcommand{\algP}{\ptxt{\textbf{P}}}
\newcommand{\algNP}{\ptxt{\textbf{NP}}}
\newcommand{\algNPC}{\ptxt{\textbf{NP-Complete}}}
\newcommand{\algNPH}{\ptxt{\textbf{NP-Hard}}}
\newcommand{\algEXP}{\ptxt{\textbf{EXP}}}
\DeclareMathOperator{\fl}{fl}

%%%%%%%%%%%%%%%%%%%%%%%%%%%%%
% Other commands
%%%%%%%%%%%%%%%%%%%%%%%%%%%%%
\newcommand{\flag}[1]{\textbf{\textcolor{red}{#1}}}
\let\uSym\u
\let\u\relax
\newcommand{\u}[1]{\underline{#1}}
\let\bSym\b
\let\b\relax
\newcommand{\b}[1]{\textbf{#1}}
\let\iSym\i
\let\i\relax
\newcommand{\i}[1]{\textit{#1}}
\let\scSym\sc
\let\sc\relax
\newcommand{\sc}[1]{\textsc{#1}}

%%%%%%%%%%%%%%%%%%%%%%%%%%%%%%%%%%%%%%%
% Make l's curvy in math environments %
%%%%%%%%%%%%%%%%%%%%%%%%%%%%%%%%%%%%%%%
\mathcode`l="8000
\begingroup
\makeatletter
\lccode`\~=`\l
\DeclareMathSymbol{\lsb@l}{\mathalpha}{letters}{`l}
\lowercase{\gdef~{\ifnum\the\mathgroup=\m@ne \ell \else \lsb@l \fi}}%
\endgroup

%%%%%%%%%%%%%%%%%%%%%%%%%
% Fix \vdots and \ddots %
%%%%%%%%%%%%%%%%%%%%%%%%%
\usepackage{letltxmacro}
\LetLtxMacro\orgvdots\vdots
\LetLtxMacro\orgddots\ddots

\makeatletter
\DeclareRobustCommand\vdots{%
	\mathpalette\@vdots{}%
}
\newcommand*{\@vdots}[2]{%
	% #1: math style
	% #2: unused
	\sbox0{$#1\cdotp\cdotp\cdotp\m@th$}%
	\sbox2{$#1.\m@th$}%
	\vbox{%
		\dimen@=\wd0 %
		\advance\dimen@ -3\ht2 %
		\kern.5\dimen@
		% remove side bearings
		\dimen@=\wd2 %
		\advance\dimen@ -\ht2 %
		\dimen2=\wd0 %
		\advance\dimen2 -\dimen@
		\vbox to \dimen2{%
			\offinterlineskip
			\copy2 \vfill\copy2 \vfill\copy2 %
		}%
	}%
}
\DeclareRobustCommand\ddots{%
	\mathinner{%
		\mathpalette\@ddots{}%
		\mkern\thinmuskip
	}%
}
\newcommand*{\@ddots}[2]{%
	% #1: math style
	% #2: unused
	\sbox0{$#1\cdotp\cdotp\cdotp\m@th$}%
	\sbox2{$#1.\m@th$}%
	\vbox{%
		\dimen@=\wd0 %
		\advance\dimen@ -3\ht2 %
		\kern.5\dimen@
		% remove side bearings
		\dimen@=\wd2 %
		\advance\dimen@ -\ht2 %
		\dimen2=\wd0 %
		\advance\dimen2 -\dimen@
		\vbox to \dimen2{%
			\offinterlineskip
			\hbox{$#1\mathpunct{.}\m@th$}%
			\vfill
			\hbox{$#1\mathpunct{\kern\wd2}\mathpunct{.}\m@th$}%
			\vfill
			\hbox{$#1\mathpunct{\kern\wd2}\mathpunct{\kern\wd2}\mathpunct{.}\m@th$}%
		}%
	}%
}
\makeatother

\newcommand{\B}{
	\begin{tikzpicture}
	\filldraw [fill=red, draw=black] (0, 0) rectangle (0.37, 0.45);
	\draw [line width=0.5mm, white ] (0.1,0.08) -- (0.1,0.38);
	\draw[line width=0.5mm, white ] (0.1, 0.35) .. controls (0.2, 0.35) and (0.4, 0.2625) .. (0.1, 0.225);
	\draw[line width=0.5mm, white ] (0.1, 0.225) .. controls (0.2, 0.225) and (0.4, 0.1625) .. (0.1, 0.1);
	\end{tikzpicture}
}

% Allow custom symbols as arrows in tikz-cd
\tikzset{
  symbol/.style={
    draw=none,
    every to/.append style={
      edge node={node [sloped, allow upside down, auto=false]{$#1$}}}
  }
}
% Example: \arrow[r,symbol=\cong]
% https://tex.stackexchange.com/questions/394154/how-to-include-inclusion-subgroup-relationship-in-tikz-cd-diagram

\author{Professor Alejandro Uribe-Ahumada\\ \small\i{Transcribed by Thomas Cohn}}
\title{Math 591 Lecure 28}
\date{11/6/20} % Can also use \today

\begin{document}
\maketitle
\setlength\RaggedRightParindent{\parindent}
\RaggedRight

\par\noindent
Recall: For $G$ a Lie group, $\mathfrak{g}=T_{e}G$ is its Lie algebra. $\forall{}A\in\mathfrak{g}$, $A^{\sharp}$ is the left-invariant vector field on $G$ determined by $A$, that is, $\forall{}g\in{}G$, $A^{\sharp}_{g}\eqdef(L_{g})_{*,e}(A)$. In particular, $A^{\sharp}_{e}=A$.\n

\defn{
	$\forall{}A\in\mathfrak{g}$, $t\in\R$, $\exp{}t\,A$ is the integral curve of $A^{\sharp}$ at $e$, at time $t$.\n
}

\defn{
	$\map{\exp:\mathfrak{g}}{G}{A}{\restr{\exp{}t\,A}{t=1}}$\n
}

\prop{
	$\exp$ is smooth.\nn
	Proof: Introduce the field $X\in\mathfrak{X}(G\times\mathfrak{g})$, defined by, $\forall{}(g,A)\in{}G\times\mathfrak{g}$,
	\[
		X_{(g,A)}=(A^{\sharp}_{g},0)\in{}T_{g}G\times{}T_{A}\mathfrak{g}=T_{(g,A)}(G\times\mathfrak{g})=T_{g}G\times{}\mathfrak{g}
	\]
	$X$ is smooth. Its flow, denoted $F$, is $F_{t}(g,A)=(F_{t}^{A}(g),A)$, where $F^{A}$ is the flow of $A^{\sharp}$. Thus, the time-$1$ map is $C^{\infty}$.
	% \[
	% 	\begin{tikzcd}[ampersand replacement=\&]
	% 		g \arrow[r, symbol=\cong]
	% 		\arrow[rrrr, bend right=15, phantom, ""{coordinate, name=Z}]
	% 		\arrow[rrrr, to path={
	% 		--(\tikztostart.south)
	% 		|-(Z)
	% 		-|(\tikztotarget.south)
	% 		--(\tikztotarget.south)
	% 		}, "\exp", swap]
	% 		\&[-15pt] \set{e}\times\mathfrak{g} \arrow[r, hook] \& G\times\mathfrak{g} \arrow[r, "F_{1}"] \& G\times\mathfrak{g} \arrow[r, "\pi_{G}"] \& G
	% 	\end{tikzcd}
	% \]
	% \[
	% 	\begin{tikzcd}[ampersand replacement=\&]
	% 		g \arrow[r, symbol=\cong]
	% 		\arrow[rr, bend right=25, phantom, ""{coordinate, name=Z1}]
	% 		\&[-15pt] \set{e}\times\mathfrak{g} \arrow[r, hook] \& G\times\mathfrak{g}
	% 		\arrow[rr, bend right=25, phantom, ""{coordinate, name=Z2}]
	% 		\arrow[r, "F_{1}"] \& G\times\mathfrak{g} \arrow[r, "\pi_{G}"] \& G
	% 		\arrow[llll, to path={
	% 			--(\tikztostart.south)
	% 			|-(Z2)
	% 			--(Z1)\tikztonodes
	% 			-|(\tikztotarget.south)
	% 			--(\tikztotarget.south)
	% 		}, "\exp", leftarrow]
	% 	\end{tikzcd}
	% \]
	\[
		\begin{tikzcd}[ampersand replacement=\&]
			g \arrow[r, symbol=\cong]
			\&[-15pt]
			\set{e}\times\mathfrak{g} \arrow[r, hook]
			\arrow[start anchor={[xshift=-1em]south}, end anchor={[xshift=1em]south}, bend right=90, phantom, ""{coordinate, name=Z1}]
			\& G\times\mathfrak{g}
			\arrow[r, "F_{1}"]
			\&
			G\times\mathfrak{g} \arrow[r, "\pi_{G}"]
			\arrow[start anchor={[xshift=-1em]south}, end anchor={[xshift=1em]south}, bend right=90, phantom, ""{coordinate, name=Z2}]
			\&
			G
			\arrow[llll, to path={
				--(\tikztostart.south)
				|-(Z2)
				--(Z1)\tikztonodes
				-|(\tikztotarget.south)
				--(\tikztotarget.south)
			}, "\exp", leftarrow, pos=0.48]
		\end{tikzcd}
	\]
	Because $F_{1}(e,A)=(\underbrace{F^{A}_{1}(e)}_{\exp(A)},A)$. We conclude that $\exp$ is smooth.\proven
}

\prop{
	With the same notation -- $A\in{}\mathfrak{g}$, $F^{A}$ the flow of $A^{\sharp}$ -- $\forall{}g\in{}G$, $F_{t}^{A}(g)=g\cdot\exp{}t\,A$ (with $\cdot$ group multiplication).\nn
	Proof: $\frac{d}{dt}(g\cdot\exp{}t\,A)=\frac{d}{dt}L_{g}(\exp{}t\,A)=(L_{g})_{*,\exp(A)}(A^{\sharp}_{\exp{}t\,A})=A^{\sharp}_{g\cdot\exp{}t\,A}$.\n
	This is precisely the ODE satisfied by $t\mapsto{}F^{A}_{t}(g)$.\proven
}

\defn{
	If $H,G$ are Lie groups, then a \u{Lie group morphism} from $H$ to $G$ is a smooth map $F:H\to{}G$ that is also a group homomorphism.\n
}

\prop{
	Let $F:H\to{}G$ be a Lie group homomorphism. Then
	\begin{enumerate}[label=\alph*), leftmargin=4\parindent]
		\item $\forall{}A\in\mathfrak{h}=T_{e}H$, $A^{\sharp}\in\mathfrak{X}(H)$ and $[\underbrace{F_{*,e}(A)}_{\in{}T_{e}G=\mathfrak{g}}]^{\sharp}\in\mathfrak{X}(G)$. These two fields are $F$-related.
		\item $\forall{}A\in\mathfrak{h}$, $F(\underbrace{\exp{}}_{\ptxt{in }H}A)=\underbrace{\exp}_{\ptxt{in }G}(F_{*,e}(A))$.
		\item $F_{*,e}:\mathfrak{h}\to\mathfrak{g}$ is a Lie algebra morphism. That is, $\forall{}A,B\in\mathfrak{h}$, $F_{*,e}(\brack{A,B})=\brack{F_{*,e}(A),F_{*,e}(B)}$ (with the appropriate Lie brackets).
	\end{enumerate}\up\n
	Proof: (b) and (c) follow directly from (a).\n
	Claim: $\forall{}h\in{}H$, the following diagram commutes:
	\[
		\begin{tikzcd}[ampersand replacement=\&]
			H \arrow[r, "F"] \arrow[d, "L_{h}", swap] \& G \arrow[d, "L_{F(h)}"]\\
			H \arrow[r, "F", swap] \& G
		\end{tikzcd}
	\]
	Check: let $k\in{}H$. Then
	\[
		(F\of{}L_{h})(k)=F(hk)=F(h)F(k)=L_{F(h)}(F(k))
	\]
	because $F$ is a group morphism. Thus, $\forall{}h$, $F_{*,h}\of(L_{h})_{*,e}=(L_{F(h)})_{*,F(h)}\of{}F_{*,e}$ by the chain rule, and commutativity at the identity). We can apply this to any $A\in\mathfrak{h}$:
	\[
		F_{*,h}(\underbrace{(L_{h})_{*,e}(A)}_{=A^{\sharp}_{h}})=(L_{F(h)})_{*,F(e)}(F_{*,e}(A))=\underbrace{(L_{F(h)})_{*,e}(F_{*,e}(A))}_{=[F_{*,e}(A)]^{\sharp_{F(h)}}}
	\]
	This \i{is} the fact that $A^{\sharp}$ and $[F_{*,e}(A)]^{\sharp}$ are $F$-related.\proven
}

\par\noindent
Now, back to subgroups.\n

\defn{
	A \u{regular} (or \u{embedded}) \u{Lie subgroup} of $G$ is a regular submanifold that is also a subgroup.\n
}

\cor{
	If $H\subseteq{}G$ is a regular subgroup, then
	\begin{enumerate}[leftmargin=4\parindent, label=(\alph*)]
		\item $\mathfrak{h}\subset\mathfrak{g}$ as a subspace is closed under $\brack{\,\cdot\,,\,\cdot\,}$ of $\mathfrak{g}$.
		\item $\forall{}A\in\mathfrak{h}$, $\underbrace{\exp}_{\ptxt{in }H}A=\underbrace{\exp}_{\ptxt{in }G}A$.
	\end{enumerate}\up\n
	Proof: Just apply the previous theorem to $F:H\hookrightarrow{}G$.\proven
}

\par\noindent
Observe: $\forall{}A,B\in\mathfrak{h}$, $\brack{A,B}_{\mathfrak{h}}=\brack{A,B}_{\mathfrak{g}}$, so $\brack{\,\cdot\,,\,\cdot\,}_{\mathfrak{h}}=\restr{\brack{\,\cdot\,,\,\cdot\,}_{\mathfrak{g}}}{\mathfrak{h}}$.\n

\defn{
	Given a Lie algebra $\mathfrak{g}$, a \u{Lie-subalgebra} is a subspace $\mathfrak{h}\subset\mathfrak{g}$ which is closed under $\brack{\,\cdot\,,\,\cdot\,}$.\n
}

\cor{
	If $H\subset{}G$ is a regular Lie subgroup, then $\mathfrak{h}=T_{e}H$ is a Lie-subalgebra of $\mathfrak{g}$.\n
}

\par\noindent
Is the converse true? I.e., given $G$ a Lie group, if $\mathfrak{h}\subset{}T_{e}G=\mathfrak{g}$ is a Lie subalgebra of $\mathfrak{g}$, is there a Lie subgroup $H\subset{}G$ such that $T_{e}H=\mathfrak{h}$?\n
Answer: Yes, if we allow immersed submanifolds.\n

\thm{
	If $\mathfrak{h}\subset{}T_{e}G$ is a Lie subalgebra, then there is a Lie group $H$ and an immersion $F:H\to{}G$ which is a group morphism and $F_{*,e}:T_{e}H\overset{\cong}{\to}\mathfrak{h}$.\n
}

\ex{
	An irrational line on the two-torus $S^{1}\times{}S^{1}=\mathbb{T}^{2}$,
	\[
		\map{F:\R}{\mathbb{T}^{2}}{t}{(e^{it},e^{i\alpha{}t})}
	\]
	for $\alpha$ irrational, is a Lie group morphism and an immersion. $\im(F_{*,e})$ is a $1$-dimensional subspace of $\mathfrak{t}^{2}\cong\R^{2}$. This is always a subalgebra.\n
}

\par\noindent
We're now done with Lie groups! Starting next week, we'll move on to differential forms.\n
Motivation: How do we integrate on manifolds and submanifolds?\n
(Think, for example, about surface integrals from calc 3.)

\end{document}